\documentclass[UTF8,fontset=windows,12pt]{ctexart}
%\setcounter{tocdepth}{1}
%\setcounter{secnumdepth}{4}
%\usepackage{ctex} %若要输入中文请取消注释或者把article改成ctexart
% This first part of the file is called the PREAMBLE. It includes
% customizations and command definitions. The preamble is everything
% between \documentclass and \begin{document}.
%\usepackage[OT1]{fontenc}
%\usepackage{fontspec}
%\setmainfont{Times New Roman}
%\setCJKmainfont{Microsoft YaHei}
\usepackage[left=1.91cm,right=1.91cm,top=2.54cm,bottom=2.54cm]{geometry}
\usepackage{amsfonts}
\usepackage{amsmath}
\usepackage{amssymb}
\usepackage[upint,lcgreekalpha]{stix}
\usepackage[type1]{garamondlibre}
%\usepackage{esint}
\usepackage{pgfplots}
\pgfplotsset{compat=1.18}
\usepackage{tikz}
\usetikzlibrary{arrows}
\usepackage{extarrows}
\usepackage[only,llbracket,rrbracket]{stmaryrd}
\usepackage{titlesec}
\usepackage{bm}
\usepackage{graphicx}
\usepackage{appendix}
\usepackage{tikz-cd}
\usepackage{tikz-3dplot}
\usetikzlibrary{arrows.meta, decorations.pathreplacing, calc, patterns}
\usepackage{float}
\usepackage{mathtools}
\usepackage{amsthm}
\usepackage{verbatim}              % better theorem environments
\usepackage{setspace}
\usepackage{enumitem}
\setenumerate[1]{itemsep=0pt,partopsep=0pt,parsep=\parskip,topsep=5pt}
\setitemize[1]{itemsep=0pt,partopsep=0pt,parsep=\parskip,topsep=5pt}
\setdescription{itemsep=0pt,partopsep=0pt,parsep=\parskip,topsep=5pt}
\usepackage{xcolor}
\usepackage[colorlinks,linkcolor=black,anchorcolor=black,citecolor=black]{hyperref}
%\usepackage{apacite}
\usepackage[numbers]{natbib}
%\usepackage{showkeys}
\usepackage{cases}
\usepackage{fancyhdr}
\usepackage[scaled=0.92]{helvet}	% set Helvetica as the sans-serif font
%\usepackage{upgreek}
%\renewcommand{\rmdefault}{ptm}		% set Times as the default text font
%\usepackage{newtxtext}
\bibliographystyle{plain}
% If AMS-LaTeX is used, it can be loaded before or after mtpro2		
% The following loads metro and defines some common MTPro options [2, 4]
%\usepackage[lite,subscriptcorrection,nofontinfo]{mtpro2}
%\pdfmapfile{+mtpro2.map}

%\usepackage{txfonts}

% various theorems, numbered by section

\makeatletter
\renewenvironment{proof}[1][\proofname]{\par
  \pushQED{\qed}%
  \normalfont \topsep6\p@\@plus6\p@\relax
  \trivlist
  \item[\hskip\labelsep
        \bfseries % 关键修改：斜体改为加粗
        #1\@addpunct{.}]\ignorespaces
}{%
  \popQED\endtrivlist\@endpefalse
}
\makeatother
\newtheoremstyle{kaiti-theorem}   % 样式名称
  {3pt}                           % 上方间距
  {3pt}                           % 下方间距
  {\kaishu\upshape}               % 正文字体：中文楷体 + 英文直立
  {}                              % 缩进量
  {\bfseries\upshape}      % 头部字体（定理名）：中文楷体粗体 + 英文直立粗体
  {.}                             % 定理名后标点
  {0.5em}                         % 定理名后间距
  {}                              % 头部说明（留空）

% 应用自定义样式
\theoremstyle{kaiti-theorem}
\newtheorem{thm}{定理}[section]
\newtheorem{nota}[thm]{记号}
\newtheorem{question}{问题}
\newtheorem*{questionn}{思考}
\newtheorem{exe}{习题}[section]
\newtheorem{assump}[thm]{假设}
\newtheorem*{thmm}{定理}
\newtheorem{defn}{定义}[section]
\newtheorem{lem}[thm]{引理}
\newtheorem*{lemm}{引理}
\newtheorem{prop}[thm]{命题}
\newtheorem*{propp}{命题}
\newtheorem*{claim}{断言}
\newtheorem{cor}[thm]{推论}
\newtheorem{conj}[thm]{猜想}
\newtheorem{rmk}{注记}[section]
\newtheorem{ex}{例}[section]
\newtheorem{pf}{证明}
\newtheorem{problem}{作业题}
\numberwithin{equation}{section}





\DeclareMathOperator{\id}{id}
\DeclareMathOperator*{\esssup}{ess\,sup}
\renewcommand{\Re}{\textbf{Re}\,}  
\renewcommand{\Im}{\textbf{Im}\,}  
\newcommand{\R}{\mathbb{R}}  
\newcommand{\FH}{\mathbb{H}}  
\newcommand{\C}{\mathbb{C}}  
\newcommand{\Z}{\mathbb{Z}}
\newcommand{\X}{\mathbb{X}}
\newcommand{\N}{\mathbb{N}}
\newcommand{\Q}{\mathbb{Q}}
\newcommand{\T}{\mathbb{T}}
\newcommand{\TT}{\mathcal{T}}
\newcommand{\TP}{\overline{\partial}{}}
\newcommand{\TJ}{\langle\TP\rangle}
\newcommand{\TL}{\overline{\Delta}{}}
\newcommand{\curl}{\text{curl }}
\newcommand{\dive}{\text{div }}
\newcommand{\bd}[1]{\mathbf{#1}}  % for bolding symbols
\newcommand{\bds}[1]{\boldsymbol{#1}}  % for bolding symbols
\newcommand{\RR}{\mathcal{R}}      % for Real numbers
\newcommand{\sh}{\mathcal{S}}  
\newcommand{\sph}{\mathbb{S}}  
\newcommand{\Om}{\Omega}
\newcommand{\OmT}{{\Omega_T}}
\newcommand{\ut}{{U_T}}
\newcommand{\vp}{\varphi}
\newcommand{\Th}{\Theta}
\newcommand{\Lam}{\Lambda}
\newcommand{\Omb}{{\overline{\Omega}}}
\newcommand{\OmbT}{{\overline{\Omega_T}}}
\newcommand{\ub}{{\overline{U}}}
\newcommand{\ubt}{{\overline{U_T}}}
\newcommand{\q}{\quad}
\newcommand{\p}{\partial}
\newcommand{\pb}{\boldsymbol{\partial}}
\newcommand{\lap}{\Delta}
\newcommand{\dd}{\mathfrak{D}}
\newcommand{\DD}{\mathcal{D}}
\newcommand{\den}{{\bar\rho}_{0}}
\newcommand{\Dt}{{\p}_{t}+v^{k}{\p}_{k}}
\newcommand{\col}[1]{\left[\begin{matrix} #1 \end{matrix} \right]}
\newcommand{\noi}{\noindent}
\newcommand{\nab}{\nabla}
\newcommand{\cnab}{\overline{\nab}}
\newcommand{\cp}{\overline{\partial}{}}
\newcommand{\dx}{\,\mathrm{d}x}
\newcommand{\dxx}{\,\mathrm{d}\bds{x}}
\newcommand{\xxi}{\boldsymbol{\xi}}
\newcommand{\eeta}{\boldsymbol{\eta}}
\newcommand{\dxxi}{\,\mathrm{d}\boldsymbol{\xi}}
\newcommand{\deeta}{\,\mathrm{d}\boldsymbol{\eta}}
\newcommand{\dxi}{\,\mathrm{d}\xi}
\newcommand{\deta}{\,\mathrm{d}\eta}
\newcommand{\dy}{\,\mathrm{d}y}
\newcommand{\dyy}{\,\mathrm{d}\bds{y}}
\newcommand{\dz}{\,\mathrm{d}z}
\newcommand{\dzz}{\,\mathrm{d}\bds{z}}
\newcommand{\dph}{\,\mathrm{d}\varphi}
\newcommand{\dt}{\,\mathrm{d}t}
\newcommand{\dr}{\,\mathrm{d}r}
\newcommand{\dtt}{\,\mathrm{d}\tau}
\newcommand{\dS}{\,\mathrm{d}S}
\newcommand{\dss}{\,\mathrm{d}\sigma}
\newcommand{\dmu}{\,\mathrm{d}\mu}
\newcommand{\dnu}{\,\mathrm{d}\nu}
\newcommand{\dV}{\,\mathrm{d}V}
\newcommand{\ds}{\,\mathrm{d}s}
\newcommand{\drr}{\,\mathrm{d}\rho}
\newcommand{\dist}{\text{dist }}
\newcommand{\vol}{\text{vol}\,}
\newcommand{\volo}{\text{vol}(\Omega)}
\newcommand{\spt}{\text{Spt}\,}
\newcommand{\Tr}{\text{Tr}\,}
\newcommand{\lee}{\langle}
\newcommand{\ree}{\rangle}
\newcommand{\xee}{\langle\xi\rangle}
\newcommand{\xxee}{\langle\boldsymbol{\xi}\rangle}
\newcommand{\xeta}{\langle\eta\rangle}
\newcommand{\xeeta}{\langle\boldsymbol{\eta}\rangle}
\newcommand{\kk}{\kappa}
\newcommand{\lam}{\lambda}
\newcommand{\io}{\int_{\Omega}}
\newcommand{\iu}{\int_{U}}
\newcommand{\is}{\int_{\Sigma}}
\newcommand{\ipo}{\int_{\partial\Omega}}
\newcommand{\ipu}{\int_{\partial U}}
\newcommand{\ig}{\int_{\Gamma}}
\newcommand{\ir}{\int_{\R}}
\newcommand{\ird}{\int_{\R^d}}
\newcommand{\irn}{\int_{\R^n}}
\newcommand{\fpi}{\frac{1}{(2\pi)^{\frac{d}{2}}}}
\newcommand{\ddt}{\frac{\mathrm{d}}{\mathrm{d}t}}
\newcommand{\ddx}{\frac{\mathrm{d}}{\mathrm{d}x}}
\newcommand{\eps}{\varepsilon}
\providecommand{\jump}[1]{\left\llbracket #1 \right\rrbracket }
\providecommand{\len}[1]{\lee #1 \ree }
\providecommand{\ino}[1]{\left\| #1 \right\| }
\providecommand{\bno}[1]{\left| #1 \right| }
\newcommand{\red}{\textcolor{red}}
\newcommand{\green}{\textcolor{green}}
\newcommand{\blue}{\textcolor{blue}}
\newcommand{\nnr}{{[n+1]}}
\newcommand{\nnn}{{[n]}}
\newcommand{\nnl}{{[n-1]}}
\newcommand{\nnll}{{[n-2]}}	
\newcommand{\mmr}{{[m+1]}}
\newcommand{\mmm}{{[m]}}
\newcommand{\mml}{{[m-1]}}
\newcommand{\mmll}{{[m-2]}}

%---------------Special variables--------------
\newcommand{\CC}{\mathfrak{C}}  
\newcommand{\NN}{\mathbf{N}}
\newcommand{\LH}{\mathcal{L}}
\newcommand{\HH}{\mathcal{H}}  
\newcommand{\EE}{\mathcal{E}}  
\newcommand{\FT}{\mathcal{F}}  
\newcommand{\fT}{\mathbf{T}}  
\newcommand{\FA}{\mathbf{A}}  
\newcommand{\MH}{\mathcal{M}}
\newcommand{\NH}{\mathcal{N}}
\newcommand{\PH}{\mathcal{P}}
\newcommand{\VH}{\mathcal{V}}
\newcommand{\XX}{\mathfrak{X}}
\newcommand{\xx}{{\bds{x}}}
\newcommand{\bx}{\mathbf{x}}
\newcommand{\by}{\mathbf{y}}
\newcommand{\bp}{\mathbf{p}}
\newcommand{\bq}{\mathbf{q}}
\newcommand{\pp}{{\bds{p}}}
\newcommand{\qq}{{\bds{q}}}
\newcommand{\uu}{\mathbf{u}}
\newcommand{\ww}{\mathbf{w}}
\newcommand{\vv}{\mathbf{v}}
\newcommand{\yy}{{\bds{y}}}
\newcommand{\zz}{{\bds{z}}}
\newcommand{\zo}{\mathbf{0}}
\newcommand{\ee}{\mathbf{e}}
\newcommand{\fa}{\mathbf{a}}
\newcommand{\fb}{\mathbf{b}}
\newcommand{\fc}{\mathbf{c}}
\newcommand{\ff}{\bds{f}}
\newcommand{\fr}{\mathbf{r}}
\newcommand{\fh}{\mathbf{h}}
\newcommand{\fm}{\mathbf{m}}
\newcommand{\fg}{\mathbf{g}}
\newcommand{\FF}{\mathbf{F}}
\newcommand{\GG}{\mathbf{G}}
\newcommand{\BB}{\mathbf{B}}



\begin{document}
\title{\LARGE{\bf  2026年春季学期现代偏微分方程作业四}}
\author{Fourier分析、Schr\"odinger方程}
\date{截止时间：2026.5.26下课前（电子版截至当天23:59:59）}

\maketitle

\begin{problem}[习题D.1.2，海森堡不确定性原理]\label{D-1:exe2}
给定点$\xx_0,\xxi_0\in \R^d$ 以及\textbf{复值}函数 $f\in\mathcal{S}(\R^d)$, 证明如下海森堡不确定性原理：
\begin{equation}
\left(\ird |(\xx-\xx_0)f(\xx)|^2\dxx\right)\left(\ird |(\xxi-\xxi_0)\hat{f}(\xxi)|^2\dxxi\right)\geqslant \frac{d^2}{4} \left(\ird |f(\xx)|^2\dxx\right)^2.
\end{equation} 这个不等式表明动量和位置不可能同时在给定的动量$\xxi_0$和给定的位置$\xx_0$附近被确定。

提示：只需证明$\xxi_0=\xx_0=0$的情况即可，否则考虑$g(\xx)=f(\xx+\xx_0)e^{-i\xx\cdot\xxi_0}$并利用命题D.1.3(2)约化到这一特殊情况。利用Plancherel恒等式可得 $|\xxi\hat{f}(\xxi)|^2=|\widehat{\nab f}(\xxi)|^2$，之后再用Plancherel恒等式和 Cauchy--Schwarz不等式证明左边$\geqslant (\ird |(\xx\cdot\nab f)f|\dxx)^2$, 最后用$\Re(\nab f)\bar{f}=\frac12 \nab(|f|^2)$, 然后分部积分一次。
\end{problem}

\begin{problem}[习题D.2.2] 若 $F\in\DD'$ 且一阶分布导数皆为零，即$\p_j F=0$ 对 $1\leqslant j\leqslant d$ 成立。证明: $F$是$\R^d$ 上的常值函数。（提示：考虑 $F*\eta_\eps$.）
\end{problem}


\begin{problem}[习题5.1.3]
设 $s_0\leqslant s\leqslant s_1$，证明：$\dot{H}^{s_0}(\R^d)\cap \dot{H}^{s_1}(\R^d)\subseteq \dot{H}^s(\R^d)$，且若 $s=(1-\theta)s_0+\theta s_1$, 则 $\|u\|_{\dot{H}^s(\R^d)}\leqslant \|u\|_{\dot{H}^{s_0}(\R^d)}^{1-\theta}\|u\|_{\dot{H}^{s_1}(\R^d)}^\theta.$ 该结论实际上对非齐次Sobolev空间也对。
\end{problem}


\begin{problem}[习题5.1.5]
设 $0<s<1,~u\in\dot{H}^s(\R^d)$. 证明：存在代表元 $u\in L_{\text{loc}}^2(\R^d)$ 且满足
\[
\int_{\R^d} \int_{\R^d} \frac{|u(\xx+\yy)-u(\xx)|^2}{|\yy|^{d+2s}}\dxx\dyy<\infty.
\]这给出了齐次Sobolev范数与Sobolev--Slobodecki\u{\i}范数之间的等价性。

提示：将 $\hat{u}$ 分解为 $\{|\xxi|\leqslant 1\}$ 部分和 $\{|\xxi|>1\}$ 部分。然后在Sobolev--Slobodecki\u{\i}范数中对 $\xx$ 变量用Plancherel恒等式。
\end{problem}

\begin{problem}[习题5.2.4]
若 $H^s(\R^d)\subset C_0(\R^d)$，其中 $C_0$ 表示在无穷远处收敛到零的连续函数，证明 $s>d/2$.

提示：用闭图像定理证明包含映射 $H^s\hookrightarrow C_0$ 是连续的，然后考虑代入Fourier变换为 $\len{\xxi}^{-2s}\chi_{|\xxi|\leqslant R}(\xxi)$ 的函数去算$s$的范围。
\end{problem}

\begin{problem}[问题5.2.1]
给定 $m\in\N^*$，设 $P(\p)=\sum\limits_{|\alpha|\leqslant m} c_\alpha\p^\alpha$ 其中 $c_\alpha\in\R$ 且存在$m$阶多重指标$\alpha$ 使得 $c_\alpha\neq 0$. 我们定义微分算子 $P(\p)$ 的\textbf{主象征 (principal symbol)} $P_m$ 为 $$P_m(\xxi):=\sum_{|\alpha|=m}c_\alpha\xxi^\alpha.$$ 我们称 $P(\p)$ 是 \textbf{$\mathbf{m}$ 阶椭圆的}，是指对任意 $\xxi\neq \bd{0}$ 都有 $P_m(\xxi)\neq 0$. 今假设 $P(\p)$ 是 $m$ 阶微分算子。
\begin{itemize}
\item [(1)] 证明：$P(\p)$ 是椭圆的当且仅当存在 $C,R>0$ 使得对 $|\xxi|\geqslant R$ 有 $|P(\xxi)|\geqslant C|\xxi|^m$.
\item [(2)] 设$P(\p)$ 是椭圆的，若 $u\in H^s(\R^d)$ 且 $P(\p) u\in H^s(\R^d)$，证明：$u\in H^{s+m}(\R^d)$.
\end{itemize}
\end{problem}

\begin{problem}[习题6.2.1]
证明非线性Schr\"odinger方程 $\bd{i}u_t+\lap u=\mu|u|^{p-1}u~(\mu=\pm 1, p\geqslant 1)$ 的质量、能量和动量守恒律。
\begin{itemize}
\item (质量守恒) $\ddt\ird |u(t,\xx)|^2\dxx=0$.
\item (能量守恒) $\ddt\ird \frac12|\nab u(t,\xx)|^2+\frac{\mu}{p+1}|u
(t,\xx)|^{p+1}\dxx=0$.
\item (动量守恒) $\ddt\Im\ird u(t,\xx)\overline{\nab u(t,\xx)}\dxx=0$.
\end{itemize}
\end{problem}

\begin{problem}[习题6.1.2+问题6.2.3]
今考虑cubic NLS解的性质。
\begin{equation}\label{6-2:cubicNLS}\tag{NLS}
\bd{i}\p_t u + \lap u = \mu |u|^2u~\text{ in }I\times\R^3,~~~u(0,\xx)=u_0(\xx).
\end{equation}其中初值$u_0\in H^1(\R^3)$, $\mu=\pm 1$, $I\subset\R$是时间区间。
\begin{itemize}
\item [(1)] 考虑使用容许对$(p,q)=(4,3)$. 令$F(u)(t,\xx)=|u|^2u(t,\xx)$，证明如下带导数的Strichartz估计
\[
\left\|\int_\R e^{i(t-\tau)\lap}F(u)(\tau,\cdot)\dtt\right\|_{L_t^\infty H_\xx^1(I\times\R^3)}\leqslant \|u\|_{L_t^4 W_\xx^{1,3}(I\times\R^3)}^3.
\]
\item [(2)] 利用(1)证明：当$\|u_0\|_{H^1(\R^3)}\ll 1$时，方程\eqref{6-2:cubicNLS}有小初值整体解  $u\in L_t^{\infty}(\R;H_\xx^1(\R^3))\cap L_t^{4}(\R;W_\xx^{1,3}(\R^3))$. 
\item [(3)] 在(2)的条件下证明散射结论：存在$u_\pm\in H^1(\R^3)$ 使得 $\lim\limits_{t\to\pm\infty}\|u(t,\cdot)-e^{it\lap}u_\pm\|_{H^1}=0$.
\item [(4)] 现在设$\mu=-1$，初值还满足$|\xx|u_0\in L^2(\R^3)$. 设$u:[0,T_*)\times\R^3\to \C$是此时\eqref{6-2:cubicNLS}在$[0,T_*)\times\R^3$上的解。定义能量$E(t)$和位力势$V(t)$如下
$$E(t):=\int_{\R^3}\frac12 |\nab u(t,\xx)|^2-\frac14 |u(t,\xx)|^4\dxx,\q V(t):=\int_{\R^3}|\xx|^2|u(t,\xx)|^2\dxx.$$ 证明：若$E(0)<0$，则$T_*<+\infty$. (提示：先写出Virial恒等式，然后证明$V''(t)\leqslant 16 E(t)$.)
\end{itemize}
\end{problem}



\end{document}
